%!TEX root = ../../main.tex
% ============================================================
% 数据表格 / 互化表
% 本文件由 main.tex 通过 \input 引入，请勿单独编译
% 重构日期: 2026-04-11
% ============================================================

\subsection{解答：分数、小数、百分数互化填表}

\vspace{0.4cm}

\noindent\textbf{7. 填表。（除不尽的保留三位小数）}

\vspace{0.8cm}

\begin{center}
\renewcommand{\arraystretch}{2.2}
\setlength{\tabcolsep}{22pt}
\begin{tabular}{|c|c|c|c|}
\hline
小\quad 数 & 0.25 & \textcolor{red}{0.667} & \textcolor{red}{1.2} \\
\hline
分\quad 数 & \textcolor{red}{$\dfrac{1}{4}$} & $\dfrac{2}{3}$ & \textcolor{red}{$\dfrac{6}{5}$} \\
\hline
百分数 & \textcolor{red}{$25\%$} & \textcolor{red}{$66.7\%$} & $120\%$ \\
\hline
\end{tabular}
\end{center}

\vspace{0.6cm}
{\small\color{gray}
\textbf{解析：}\\
第一列：已知小数 $0.25$。转化为分数为 $\dfrac{25}{100} = \dfrac{1}{4}$，化为百分数将小数点向右移两位，即 $25\%$。\\
第二列：已知分数 $\dfrac{2}{3}$。转化为小数为 $2 \div 3 \approx 0.667$（保留三位小数），化为百分数为 $66.7\%$。\\
第三列：已知百分数 $120\%$。转化为分数为 $\dfrac{120}{100} = \dfrac{6}{5}$，化为小数为也就是去掉百分号并将小数点向左移两位，即 $1.2$。
}

\vspace{0.8cm}

%% ==============================
%% 画线段图解三角形内角和
%% ==============================
\newpage
